\begin{abstract}
\noindent
English football spends heavily on finding young players, yet it is not known whether the players who reach the top two divisions are born across England in proportion to births or in particular places. We ask where English-born men who played in the Premier League or the Championship were born, compared with where boys were born. From Wikidata and Wikipedia we identify \nStudy{} men born 1975--2005 with evidence of a league appearance at one of \nClubsTopTwo{} top-two clubs since 1992; \nEnglish{} were born in England and \nMapped{} can be placed on a map. Expected counts come from \birthsTotal{} million male births estimated from ONS mid-year populations. Players are \regionRatio{} times as common per birth in London (\londonRate{} per 10\,000 male births) as in the South West (\swRate{}), and \nRegionsSig{} of 9 regions differ from England at $p<0.05$. Within regions, birthplaces are more clustered than births at distances up to 60\,km (inhomogeneous $L$-function, DCLF $p=\lDCLFall{}$; births-only null $p=\lNullAll{}$). Across Travel to Work Areas (TTWAs) the Empirical Bayes Moran's $I$ is \moranTTWA{} ($p=\moranTTWAp{}$), with high values in a belt from Merseyside and Lancashire across to South Yorkshire and low values along the south coast, but only one area (\giTopName{}, $q=\giTopQ{}$) comes close to passing a false discovery rate correction. In a negative binomial model with region effects, each extra 10\,km to the nearest professional club lowers the rate by \distPct{}\% (rate ratio \distRR{}, 95\% CI \distLo{}--\distHi{}, $p\,\distP{}$), while deprivation and urban share have no effect once region is included. For the \nAcadPlayers{} players born 1996--2005, who were of academy age after the 2012 academy reform, distance to any club matters (rate ratio \clubYoungRR{}) and distance to an elite Category 1 academy does not (\acadRR{}, $p=\acadP{}$). Of all \nStudy{} study players, \shareOutsideEngland\% were born outside England, mostly in Western Europe and the other UK nations. Of the stated hypotheses, H0 (births alone explain the pattern) is rejected, H1 (individual hotspots and coldspots) is supported for the overall pattern but only weakly for named areas, and H2 (association with deprivation, urbanisation and proximity to clubs and academies) is supported for club proximity only. The results describe areas, not individuals.
\end{abstract}
